\documentclass[11pt]{article}
\usepackage[utf8]{inputenc}
\usepackage{amsmath, amssymb, amsthm}
\usepackage[margin=1in]{geometry}

\theoremstyle{definition}
\newtheorem{definition}{Definition}
\newtheorem{theorem}{Theorem}

\theoremstyle{remark}
\newtheorem{remark}{Remark}

\DeclareMathOperator{\tr}{tr}

\title{Lecture 1: SPDE --- Introduction}
\author{}
\date{16.10.2025}
\author{Tien Anh Vu}
\begin{document}
\maketitle

%% =============================
%% Example: Dean--Kawasaki equation
%% =============================
\section*{Introduction}

\subsection*{Example: Dean--Kawasaki equation}

% Reference: Dean (1996), Kawasaki (1998).
% This is a motivating example for SPDEs arising from
% interacting particle systems.
This is a motivating example for SPDEs arising from interacting particle systems.

\noindent
\textbf{(Dean 1996, Kawasaki 1998)}

Consider independent Brownian motions $B^1, B^2, B^3, \ldots$
(independent BMs), and a smooth function $f \in C_b^2(\mathbb{R})$, i.e., twice differentiable with bounded derivatives and bounded values.

% We take a test function f in C_b^2(R), i.e., twice differentiable
% with bounded derivatives and bounded values.


%% =============================
%% Empirical average
%% =============================
\paragraph{Empirical average.}
% The empirical measure of N particles, each performing
% independent BM starting from x_0, is averaged against f.
The empirical measure of $N$ particles, each performing independent BM starting from $x_0$, is averaged against $f$.

Then for $N$ particles:
\[
    \eta_t^N(f)(\omega)
    \;:=\;
    \frac{1}{N}\sum_{i=1}^{N}
    f\!\bigl(B_t^i + x_0\bigr)(\omega)
\]
% This is the empirical (sample) average of f evaluated
% at the particle positions B_t^i + x_0.
This is the empirical (sample) average of $f$ evaluated at the particle positions $B_t^i + x_0$.

\begin{remark}
We can rewrite this as an integral against the empirical measure:
\[
    \eta_t^N(f)(\omega)
    \;=\;
    \int f\,d\!\left(
        \frac{1}{N}\sum_{i=1}^{N}
        \delta_{B_t^i + x_0}
    \right)(\omega).
\]
\end{remark}

%% =============================
%% SLLN: convergence as N -> infinity
%% =============================
\paragraph{Law of large numbers.}
% By the strong law of large numbers (SLLN),
% the empirical average converges P-a.s. as N -> infinity.
By the SLLN, the empirical average converges $P$-a.s.\ as $N \to \infty$.

By the \textbf{strong law of large numbers (SLLN)},
as $N \to \infty$:
\[
    \eta_t^N(f)
    \;\xrightarrow{P\text{-a.s.}}\;
    E\!\bigl(f(B_t + x_0)\bigr).
\]

% The distribution of B_t + x_0 is N(x_0, t), so the
% expectation is an integral against the Gaussian density.
The distribution of $B_t + x_0$ is $\mathcal{N}(x_0, t)$, so the expectation is an integral against the Gaussian density.

Since $B_t + x_0 \sim \mathcal{N}(x_0,\, t)$,
the right-hand side equals:
\[
    E\!\bigl(f(B_t + x_0)\bigr)
    \;=\;
    \int_{\mathbb{R}} f(x)\,p(t, x)\,dx,
\]
where $p(t, x)$ is the \textbf{Gaussian heat kernel}:
\[
    p(t, x)
    \;=\;
    \frac{1}{\sqrt{2\pi t}}\;
    e^{-\frac{(x - x_0)^2}{2t}}.
\]

\begin{remark}
The convergence $\eta_t^N(f) \xrightarrow{P\text{-a.s.}} E\!\bigl(f(B_t + x_0)\bigr)$ holds for each fixed~$t$ and each test function~$f$.  The right-hand side is deterministic: all randomness has been averaged out.  The Gaussian density $p(t,x)$ arises because the sum of $N$ i.i.d.\ displacements (each distributed as $B_t$) concentrates around the common mean, and each particle's position has the distribution $\mathcal{N}(x_0, t)$.
\end{remark}

%% =============================
%% Connection to the heat equation (PDE)
%% =============================
\subsection*{Connection to the heat equation}

% The density p(t,x) solves the heat equation.
% This is the classical connection between Brownian motion
% and parabolic PDEs.
The density $p(t,x)$ solves the heat equation. This is the classical connection between Brownian motion and parabolic PDEs.

Then $p(t, x)$ solves the following \textbf{PDE} (heat equation):
\[
    \boxed{
    \begin{cases}
        \displaystyle
        \partial_t\, p(t, x)
        \;=\;
        \tfrac{1}{2}\;\partial_{xx}\, p(t, x),
        & t > 0,\; x \in \mathbb{R}, \\[10pt]
        p(t, x)\,dx
        \;\longrightarrow\;
        \delta_{x_0}(dx),
        & \text{(initial condition)}.
    \end{cases}
    }
\]

\begin{remark}
This is the \textbf{heat equation}.
The initial condition says that as $t \to 0^+$,
the measure $p(t, x)\,dx$ converges (in the sense of
distributions) to the Dirac delta $\delta_{x_0}$.
% Physically: heat concentrated at x_0 at time 0
% diffuses according to the heat equation.
% The solution p(t,x) is the fundamental solution
% (Green's function) of the heat operator.
Physically: heat concentrated at $x_0$ at time $0$
diffuses according to the heat equation.
The solution $p(t,x)$ is the fundamental solution
(Green's function) of the heat operator.
\end{remark}

\begin{remark}
Intuitively, the heat equation describes diffusion:
an initial concentration of ``heat'' at $x_0$ spreads
out over time, with the profile given by the Gaussian
$p(t,x)$.  The solution has \textbf{negative curvature}
--- the temperature decreases over time at the peak
and spreads to the sides.
\end{remark}

%% =============================
%% From microscopic to macroscopic: BM -> PDE
%% =============================
\subsection*{From microscopic to macroscopic description}

% BM (microscopic description, goes on levels) -> macroscopic description (PDE).

Brownian motion provides the \textbf{microscopic description}
(individual particles).
The PDE (heat equation) provides the \textbf{macroscopic description}
(density of particles).

\begin{remark}
This passage from \emph{microscopic} to \emph{macroscopic} is a
central theme in mathematical physics and probability theory.
At the microscopic level, we track each particle individually
(a system of $N$ SDEs); at the macroscopic level, we describe
the \emph{collective behaviour} via a single PDE for the density.
The law of large numbers is the mathematical mechanism that
justifies this passage: random fluctuations average out as
$N \to \infty$, and only the deterministic drift survives.
This is sometimes called a \textbf{mean-field limit} or
\textbf{hydrodynamic limit}.
\end{remark}

%% =============================
%% Fluctuations around LLN: Itô's formula
%% =============================
\subsection*{Fluctuations around LLN --- Itô's formula}

% We apply Itô's formula to the empirical average
% to understand the fluctuations around the law of large numbers limit.

Apply Itô's formula to the empirical average:
\begin{align*}
    d\!\left(
        \frac{1}{N}\sum_{i=1}^{N}
        f\!\bigl(B_t^i + x_0\bigr)
    \right)
    &\;=\;
    \frac{1}{N}\sum_{i=1}^{N}
    df\!\bigl(B_t^i + x_0\bigr) \\[8pt]
    &\;=\;
    \underbrace{
        \frac{1}{N}\sum_{i=1}^{N}
        f'\!\bigl(B_t^i + x_0\bigr)\,dB_t^i
    }_{\text{Itô integral}}
    \;+\;
    \underbrace{
        \frac{1}{N}\sum_{i=1}^{N}
        \frac{1}{2}\,f''\!\bigl(B_t^i + x_0\bigr)\,dt
    }_{\text{correction term}}.
\end{align*}

% The first term is the stochastic (Itô) integral part.
% The second term is the drift/correction from Itô's formula,
% which equals eta_t^N(1/2 f'').

We can identify the two terms:
\begin{itemize}
    \item The Itô integral part defines a \textbf{local martingale}
          $M_t^N(f)$.
    \item The correction term equals
          $\eta_t^N\!\bigl(\tfrac{1}{2}f''\bigr)\,dt$.
\end{itemize}

So in shorthand:
\[
    d\,\eta_t^N(f)
    \;=\;
    dM_t^N(f)
    \;+\;
    \eta_t^N\!\bigl(\tfrac{1}{2}f''\bigr)\,dt.
\]

% M_t^N(f) is a local martingale.
% As N -> infinity, M_t^N(f) -> 0, recovering the PDE.
% This shows the asymptotic connection between Itô integrals
% and the heat equation.

This is the \textbf{asymptotic of the Itô integral}:
as $N \to \infty$, $M_t^N(f) \to 0$, and we recover the PDE.

\begin{remark}
The decomposition $d\,\eta_t^N(f) = dM_t^N(f) + \eta_t^N(\tfrac{1}{2}f'')\,dt$
is a \textbf{semimartingale decomposition}: every semimartingale can
be written as a local martingale plus a finite-variation (drift) part.
The martingale part $M_t^N(f)$ captures the \emph{randomness}
(fluctuations due to the individual Brownian motions), while the
drift part $\eta_t^N(\tfrac{1}{2}f'')\,dt$ captures the
\emph{deterministic trend} (the diffusion described by the heat
equation).  As $N \to \infty$, the noise averages out and only the
drift remains~--- this is the law of large numbers at the level of
processes.
\end{remark}

\begin{remark}
The factor $\tfrac{1}{2}f''$ appearing in the drift is the
\textbf{generator of Brownian motion}.  In general, if $X_t$
is a diffusion with generator $\mathcal{A}$, then It\^o's formula
gives $d\,f(X_t) = \mathcal{A}f(X_t)\,dt + \text{(martingale)}$.
For standard BM, $\mathcal{A} = \tfrac{1}{2}\partial_{xx}$,
hence the $\tfrac{1}{2}f''$ term.
\end{remark}

%% =============================
%% Quadratic variation
%% =============================
\paragraph{Quadratic variation.}

% We compute the quadratic variation of the martingale M^N(f)
% using the Itô isometry.

The quadratic variation is computed using the \textbf{Itô isometry}:
\begin{align*}
    E\!\Bigl(\bigl(M_t^N(f)\bigr)^2\Bigr)
    &\;=\;
    \frac{1}{N^2}\sum_{i=1}^{N}
    E\!\left(
        \left(\int_0^t f'\!\bigl(B_s^i + x_0\bigr)\,dB_s^i
        \right)^{\!2}
    \right) \\[8pt]
    %% By Itô isometry: E((int g dB)^2) = E(int g^2 ds)
    &\;=\;
    \frac{1}{N^2}\sum_{i=1}^{N}
    E\!\left(
        \int_0^t
        f'\!\bigl(B_s^i + x_0\bigr)^2\,ds
    \right) \\[8pt]
    &\;=\;
    \frac{1}{N}\;
    E\!\left(
        \int_0^t
        f'\!\bigl(B_s + x_0\bigr)^2\,ds
    \right).
\end{align*}

% In the last step we used that the B^i are i.i.d.,
% so all N terms in the sum are equal, giving a factor 1/N.

In the last step, we used that the $B^i$ are i.i.d.,
so all $N$ terms are equal, giving a factor $1/N$.

%% =============================
%% Quadratic variation via empirical measure
%% =============================
\paragraph{Quadratic variation process.}

% The predictable quadratic variation <M^N(f)>_t
% can be written in terms of the empirical measure.

We have
\[
    \bigl\langle M^N(f) \bigr\rangle_t
    \;=\;
    \frac{1}{N^2}\sum_{i=1}^{N}
    \int_0^t f'\!\bigl(B_s^i + x_0\bigr)^2\,ds
    \;=\;
    \frac{1}{N}
    \int_0^t \eta_s^N\!\bigl((f')^2\bigr)\,ds
    \;=\;
    O\!\left(\frac{1}{N}\right).
\]

% The quadratic variation is O(1/N), which means the
% martingale part vanishes as N -> infinity.
% This confirms the law of large numbers: the fluctuations
% around the deterministic limit are of order 1/sqrt(N).

The quadratic variation is $O(1/N)$, confirming that
the martingale part vanishes as $N \to \infty$.
The fluctuations around the deterministic limit
are of order $1/\sqrt{N}$.

\begin{remark}
The \textbf{It\^o isometry} states that for a (predictable)
integrand $\phi_s$ and a Brownian motion $B_s$:
\[
    E\!\left[\left(\int_0^t \phi_s\,dB_s\right)^{\!2}\right]
    \;=\;
    E\!\left[\int_0^t \phi_s^2\,ds\right].
\]
It converts a stochastic quantity (second moment of the It\^o
integral) into a deterministic one (expected integral of the
integrand squared).  This is why the cross terms
$i \neq j$ vanish in the computation: the Brownian motions
$B^i$ and $B^j$ are \emph{independent}, so
$E\!\bigl(\int \phi\,dB^i \cdot \int \psi\,dB^j\bigr) = 0$.
\end{remark}

\begin{remark}
The quadratic variation being $O(1/N)$ means that
$E\!\bigl((M_t^N)^2\bigr) \to 0$, which implies
$M_t^N \to 0$ in $L^2(\Omega)$ (and hence in probability).
This is a \emph{quantitative} version of the law of large numbers:
the variance of the empirical average decreases at rate $1/N$
(equivalently, the standard deviation decreases at rate $1/\sqrt{N}$).
\end{remark}

%% =============================
%% Integral form of the equation
%% =============================
\subsection*{Integral form}

% Recall ($\star$): d eta_t^N(f) = dM_t^N(f) + eta_t^N(1/2 f'') dt.
% As N -> infinity, the martingale part vanishes (-> 0).

Recall from $(\star)$:
\[
    d\,\eta_t^N(f)
    \;=\;
    dM_t^N(f)
    \;+\;
    \eta_t^N\!\bigl(\tfrac{1}{2}f''\bigr)\,dt.
\]

As $N \to \infty$, the martingale part $M_t^N(f) \to 0$
(since its quadratic variation is $O(1/N)$),
so only the drift term survives.
Integrating in time:
\[
    \eta_t^N(f)
    \;=\;
    \eta_0^N(f)
    \;+\;
    \int_0^t \eta_s^N\!\bigl(\tfrac{1}{2}f''\bigr)\,ds.
\]

In the limit $N \to \infty$, this becomes the
\textbf{weak (distributional) form} of the heat equation:
for every test function $f \in C_b^2(\mathbb{R})$,
the limiting density $\eta_t$ satisfies
$\eta_t(f) = \eta_0(f) + \int_0^t \eta_s(\tfrac{1}{2}f'')\,ds$.

\begin{remark}
The integral form $\eta_t^N(f) = \eta_0^N(f) + M_t^N(f) + \int_0^t \eta_s^N(\tfrac{1}{2}f'')\,ds$ is a \textbf{mild formulation}: rather than requiring differentiability of $\eta_t^N$ in time, the equation is expressed as an integral identity.  This is advantageous because the empirical measure $\eta_t^N$ is a random sum of Dirac masses, which is not differentiable in the classical sense.  The passage to the limit $N \to \infty$ is performed at the level of this integral identity, which is natural for probabilistic arguments where limits are taken in distribution or in $L^p$, rather than pointwise for derivatives.
\end{remark}

%% =============================
%% Density form: eta_s^N -> density
%% =============================
\paragraph{Density form.}

% The empirical measure eta_s^N has a density eta_s^N(x).
% We assume we can pass to the limit to get a density eta_s(x).

The empirical measure $\eta_s^N$ is a sum of Dirac masses.
Its density (in the distributional sense) is:
Define the empirical density:
\[
    \eta_s^N(x)
    \;:=\;
    \frac{1}{N}\sum_{i=1}^{N}
    \delta_{B_s^i + x_0}(x).
\]

Assuming we can pass to the limit (i.e., replace
$\eta_s^N$ by a smooth density $\eta_s$), we use
\textbf{integration by parts} (twice) to move the
derivatives from $f$ onto $\eta$:
\[
    \eta_t^N\!\bigl(\tfrac{1}{2}f''\bigr)
    \;=\;
    \int f''(x)\;\tfrac{1}{2}\,\eta_t^N(x)\,dx
    \;=\;
    \int f(x)\;\tfrac{1}{2}\,\partial_{xx}\,\eta_t^N(x)\,dx,
\]
where the boundary terms vanish because $f \in C_b^2$.
This gives
\[
    \eta_t^N\!\bigl(\tfrac{1}{2}f''\bigr)\,dt
    \;=\;
    \tfrac{1}{2}\,\partial_{xx}\,\eta_t^N(f)\,dt.
\]

This is precisely the \textbf{weak form of the heat equation}:
the identity $\int f \cdot \tfrac{1}{2}f''\,\eta\,dx
= \int f \cdot \tfrac{1}{2}\partial_{xx}\eta\,dx$
holds after integration by parts, confirming that the
limiting density solves $\partial_t \eta = \tfrac{1}{2}\partial_{xx}\eta$.

\begin{remark}
\textbf{Integration by parts} is the key tool that converts the
\emph{weak formulation} (derivatives on the test function~$f$) into
the \emph{strong formulation} (derivatives on the density~$\eta$).
In the weak form, one only requires $f$ to be smooth; the density
$\eta$ can be rough.  This is essential because the empirical
measure $\eta_s^N$ is a sum of Dirac masses (not differentiable),
so the weak formulation is the natural one at finite $N$.  Only
in the limit $N \to \infty$, when $\eta$ becomes smooth, can we
pass to the strong (classical) PDE form.
\end{remark}

\begin{remark}
The weak form $\eta_t(f) = \eta_0(f) + \int_0^t
\eta_s(\tfrac{1}{2}f'')\,ds$ for all test functions $f$ is
the \textbf{mild/weak solution concept} for the heat equation.
It says: ``test the density against any smooth function $f$ and
check that the resulting real-valued function of time satisfies
an ODE.''  This avoids requiring pointwise differentiability of
$\eta$ and is the standard notion of solution for PDEs with
irregular data.
\end{remark}

%% =============================
%% Fluctuations around SLLN
%% =============================
\subsection*{Fluctuations around SLLN}

% The martingale part captures the fluctuations.
% The 1st and 2nd moments coincide in the limit.

While the martingale $M_t^N(f)$ vanishes in the limit,
its \textbf{fluctuations} (after rescaling by $\sqrt{N}$)
converge to a non-trivial limit.
The martingale part is:
\[
    M_t^N(f)
    \;=\;
    \frac{1}{N}\sum_{i=1}^{N}
    \int_0^t f'\!\bigl(B_s^i + x_0\bigr)\,dB_s^i.
\]

We can rewrite this by splitting $\frac{1}{N} = \frac{1}{\sqrt{N}} \cdot \frac{1}{\sqrt{N}}$:
\[
    M_t^N(f)
    \;=\;
    \frac{1}{\sqrt{N}}
    \cdot
    \frac{1}{\sqrt{N}}\sum_{i=1}^{N}
    \int_0^t f'\!\bigl(B_s^i + x_0\bigr)\,dB_s^i.
\]

\begin{remark}
The splitting $\frac{1}{N} = \frac{1}{\sqrt{N}} \cdot \frac{1}{\sqrt{N}}$ is the key step that reveals the \textbf{central limit theorem scaling}.  The inner factor $\frac{1}{\sqrt{N}}\sum_{i=1}^N$ is the CLT-normalised sum, which converges to a non-trivial Gaussian limit.  The outer factor $\frac{1}{\sqrt{N}}$ then sends the whole expression to zero, consistently with the LLN.  Removing the outer $\frac{1}{\sqrt{N}}$ (i.e., studying $\sqrt{N}\,M_t^N$) isolates the fluctuations.
\end{remark}

The key observation is that the 1st and 2nd moments
of the rescaled sum coincide in the limit
(by the law of large numbers for the quadratic variation).

%% =============================
%% Limiting form: where eta_t(x) dx = p(t,x) dx
%% =============================
\paragraph{Limiting density.}

In the limit, $\eta_t(x)\,dx = p(t,x)\,dx$
(the Gaussian heat kernel), and the martingale part,
after rescaling by $\sqrt{N}$, converges to a
\textbf{stochastic integral} involving the
square root of the limiting density:
\[
    M_t^N(f)
    \;\approx\;
    \frac{1}{\sqrt{N}}
    \int_0^t \!\!\int
    f'(x)\,\sqrt{\eta_s^N(x)}\;
    d\xi(s, x),
\]
where $d\xi(s,x)$ denotes \textbf{space-time white noise}
(informally: one independent Brownian motion per spatial location $x$).

The appearance of $\sqrt{\eta_s^N(x)}$ is natural:
each particle at position $x$ contributes an independent
Brownian increment, and the ``density of particles'' near $x$
is $\eta_s^N(x)$, so by the central limit theorem the
fluctuations scale like $\sqrt{\eta_s^N(x)}$.

\begin{remark}
The factor $\sqrt{\eta_s^N(x)}$ can be understood from two complementary viewpoints.
\textbf{Probabilistically:} Near position $x$, there are approximately $N\,\eta_s^N(x)\,\Delta x$ particles in a small interval of length $\Delta x$.  Each contributes an independent noise of order $1$, so by the CLT the total fluctuation is of order $\sqrt{N\,\eta_s^N(x)\,\Delta x}$.  After dividing by $N$ (to form the empirical density), the noise per unit length is proportional to $\sqrt{\eta_s^N(x)/N}$.
\textbf{Analytically:} The $\sqrt{\eta}$ factor is dictated by the \emph{quadratic variation}: the It\^o isometry forces the noise coefficient to be the square root of the integrand appearing in the quadratic variation $\langle M^N(f)\rangle_t$, which is $\eta_s^N\bigl((f')^2\bigr)$.
\end{remark}

The quadratic variation confirms this:
\[
    \bigl\langle M^N(f)\bigr\rangle_t
    \;=\;
    \frac{1}{N}
    \int_0^t \!\!\int
    \bigl(f'(x)\bigr)^2\,\eta_s^N(x)\,dx\,ds,
\]
which is exactly the Itô isometry applied to the
stochastic integral representation above.

%% =============================
%% CLT for the rescaled martingale
%% =============================
\subsection*{CLT for the rescaled martingale}

% The rescaled martingale sqrt(N) M_t^N(f) converges
% to an Itô integral w.r.t. a single BM.

The rescaled martingale converges:
\[
    \sqrt{N}\;M_t^N(f)
    \;\longrightarrow\;
    \int_0^t \sqrt{\eta_s\!\bigl((f')^2\bigr)}\;dB_s
    \qquad \text{(Itô integral w.r.t.\ BM)}.
\]

% For each test function f, the limiting random variable
% is Gaussian by the martingale CLT.

For each test function $f$, the limit is Gaussian:
\[
    f \;\longmapsto\;
    \mathcal{N}\!\left(0,\;
    \int_0^t \eta_s\!\bigl((f')^2\bigr)\,ds
    \right).
\]

\begin{remark}
This is the \textbf{martingale central limit theorem} (martingale CLT)
in action.  The classical CLT says that
$\frac{1}{\sqrt{N}}\sum_{i=1}^N X_i \to \mathcal{N}(0, \sigma^2)$
for i.i.d.\ random variables.  Here, the summands are
\emph{stochastic integrals} $\int_0^t f'(B_s^i + x_0)\,dB_s^i$,
which are themselves random variables (in fact, martingales).
The martingale CLT guarantees that their normalised sum converges
to a Gaussian, with variance determined by the
\textbf{predictable quadratic variation}.
\end{remark}

\begin{remark}
The variance $\int_0^t \eta_s\!\bigl((f')^2\bigr)\,ds$ of the
limiting Gaussian depends on the \emph{test function $f$} through
$(f')^2$ and on the \emph{limiting density $\eta_s$}.
This means that the fluctuations are spatially
\emph{heterogeneous}: regions with higher particle density
$\eta_s(x)$ contribute more variance.  This is the
infinite-dimensional analogue of the classical result that
the variance of a sample mean depends on the population variance.
\end{remark}

%% =============================
%% The approximate SDE
%% =============================
\paragraph{Leading to the SDE.}

% Combining the drift and martingale parts with the
% Itô isometry identification gives the approximate SDE.

Combining the drift and the martingale part
(using the Itô isometry):
\[
    d\,\eta_t^N(f)
    \;\approx\;
    \eta_t^N\!\bigl(\tfrac{1}{2}f''\bigr)\,dt
    \;+\;
    \frac{1}{\sqrt{N}}\;
    \sqrt{\eta_t^N\!\bigl((f')^2\bigr)}\;dB_t.
\]

\begin{remark}
This approximate SDE shows the two competing effects:
\begin{itemize}
    \item The \textbf{drift} $\eta_t^N(\tfrac{1}{2}f'')\,dt$ drives
          the system towards the deterministic heat equation.
    \item The \textbf{noise} of order $1/\sqrt{N}$ captures the
          random fluctuations due to finite particle number.
\end{itemize}
As $N \to \infty$, the noise term vanishes and we recover the
deterministic PDE.  For large but finite $N$, this equation
quantifies how the empirical measure deviates from the
hydrodynamic limit.
\end{remark}

%% =============================
%% Space-time white noise representation
%% =============================
\subsection*{Space-time white noise representation}

% For s > 0, eta_s^N has a density eta_s^N(x).
% Introduce space-time white noise xi(t,x).

For $s > 0$, we can write $\eta_s^N(dx) = \eta_s^N(x)\,dx$,
and using \textbf{space-time white noise} $\xi(t,x)$,
i.e.\
\[
    \mathbb{E}\!\bigl(\xi(t,x)\,\xi(s,y)\bigr)
    \;=\;
    \delta(t - s)\;\delta(x - y)
    \qquad \text{(uncorrelated)},
\]
we can write the following chain of equalities.

% Rewriting the quadratic variation as an L^2 norm
% and identifying via the Itô isometry for space-time white noise.

The $L^2$ norm squared of the martingale
(i.e.\ the quadratic variation) can be rewritten as:
\begin{align*}
    \frac{1}{N}\;
    \mathbb{E}\!\left(
        \int_0^t \eta_s^N\!\bigl((f')^2\bigr)\,ds
    \right)
    &\;=\;
    \frac{1}{N}\;
    \mathbb{E}\!\left(
        \int_0^t \!\!\int_{\mathbb{R}}
        (f')^2(x)\;\eta_s^N(x)\,dx\,ds
    \right) \\[10pt]
    &\;=\;
    \frac{1}{N}\;
    \mathbb{E}\!\left[
        \left(
            \int_0^t \!\!\int_{\mathbb{R}}
            f'(x)\;\sqrt{\eta_s^N(x)}\;
            d\xi(s, x)
        \right)^{\!2}
    \right].
\end{align*}

% The last equality is the Itô isometry for the
% stochastic integral against space-time white noise:
% E[(int int g d xi)^2] = int int g^2 ds dx.

The last equality follows from the \textbf{Itô isometry}
for the stochastic integral against space-time white noise:
if $g(s,x) = f'(x)\,\sqrt{\eta_s^N(x)}$, then
$g^2 = (f')^2(x)\,\eta_s^N(x)$, confirming the identity.

\begin{remark}
\textbf{Space-time white noise} $\xi(t,x)$ is the
infinite-dimensional analogue of standard white noise.
Formally, $\xi$ is a generalised Gaussian random field
indexed by $(t,x) \in [0,T] \times \mathbb{R}$ with
covariance
$E\!\bigl(\xi(t,x)\,\xi(s,y)\bigr) = \delta(t-s)\,\delta(x-y)$.
Rigorously, it is defined as a random distribution:
for each test function $\varphi \in L^2([0,T] \times \mathbb{R})$,
the pairing $\xi(\varphi) := \int\!\!\int \varphi(t,x)\,\xi(t,x)\,dt\,dx$
is a centred Gaussian random variable with variance
$\|\varphi\|_{L^2}^2$.  This is the natural noise driving SPDEs,
just as standard BM is the natural noise driving SDEs.
\end{remark}

\begin{remark}
The passage from the sum of $N$ independent martingales to a
single stochastic integral against space-time white noise
is the key conceptual step.  Heuristically, each particle at
position $x$ contributes an independent ``Brownian kick'';
as $N \to \infty$, these individual kicks merge into a
continuum of independent noises parametrised by $x$, which is
precisely what space-time white noise models.
\end{remark}

%% =============================
%% Dean--Kawasaki SPDE
%% =============================
\subsection*{The Dean--Kawasaki SPDE}

% This representation leads to the Dean-Kawasaki SPDE,
% which is ill-posed (no general solution exists).

This leads to the \textbf{Dean--Kawasaki SPDE}:
\[
    \boxed{
        \partial_t\,\eta_t^N
        \;=\;
        \tfrac{1}{2}\,\partial_{xx}\,\eta_t^N(x)
        \;+\;
        \frac{1}{\sqrt{N}}\;\partial_x\!\left(
            \sqrt{\eta_t^N(x)}\;\xi(t, x)
        \right).
    }
\]

\begin{remark}
This SPDE is \textbf{ill-posed}: no general (function-valued) solution exists.
The noise term involves both a square root of the solution and a spatial
derivative of the space-time white noise, creating severe irregularity.
\end{remark}

\begin{remark}
The Dean--Kawasaki SPDE has the structure of a
\textbf{stochastic conservation law}: the noise enters through
a spatial derivative $\partial_x(\cdots)$, which ensures that
the total mass $\int \eta_t^N(x)\,dx$ is formally conserved
(the noise only \emph{redistributes} mass, it does not create
or destroy it).  This is physically natural since the number
of particles is fixed at $N$.
\end{remark}

\begin{remark}
Despite its ill-posedness as a classical SPDE, the Dean--Kawasaki
equation is extremely useful as a \textbf{formal identity}: it
correctly encodes both the law-of-large-numbers limit
(heat equation) and the fluctuations (the $1/\sqrt{N}$ noise).
Recent mathematical work has given rigorous meaning to this
equation using \textbf{measure-valued solutions} or
\textbf{renormalisation} techniques.
\end{remark}

%% =============================
%% OU process preview
%% =============================
\paragraph{OU process.}

% Preview of the Ornstein--Uhlenbeck process:
% a linear SDE that will be studied later.

As a simpler model, the \textbf{Ornstein--Uhlenbeck (OU) process}
is the solution of the linear SDE:
\[
    dX_t \;=\; A\,X_t\,dt \;+\; C\,dW_t.
\]

\begin{remark}
The OU process is the simplest non-trivial example of a
\textbf{linear SDE}.  Compared to the Dean--Kawasaki equation,
it is well-posed and explicitly solvable:
\[
    X_t \;=\; e^{At}\,X_0
    \;+\;\int_0^t e^{A(t-s)}\,C\;dW_s.
\]
Here $A$ is a (possibly unbounded) linear operator generating
mean reversion (if $A < 0$), and $C$ controls the noise intensity.
The OU process will serve as the \textbf{linear prototype}
for the general theory of SPDEs on Hilbert spaces that follows
in Chapter~1.
\end{remark}

\begin{remark}
In finite dimensions ($H = \mathbb{R}^d$), the OU process
has a Gaussian invariant measure $\mathcal{N}(0, Q_\infty)$
provided $A$ is stable (all eigenvalues have negative real part),
where $Q_\infty$ solves the \textbf{Lyapunov equation}
$A Q_\infty + Q_\infty A^* + C C^* = 0$.  In infinite dimensions,
the analogous condition requires $Q_\infty$ to be trace-class,
which imposes a \emph{compatibility condition} between the
operators $A$ and $C$.
\end{remark}

%% =============================================================
%% Chapter 1: Stochastic Analysis on Hilbert Spaces
%% =============================================================
\newpage
\section*{Chapter 1: Stochastic Analysis on Hilbert Spaces}

%% =============================
%% Hilbert spaces
%% =============================
\subsection*{Hilbert spaces}

% (L^2(M), <.,.>) is a complete normed vector space.

Recall that $\bigl(L^2(\mu),\,\langle\cdot,\cdot\rangle\bigr)$ is a
\textbf{complete normed vector space} (Hilbert space).

The inner product is given by,
\[
    \langle f, g \rangle = \int f\,g\;d\mu.
\]

% A linear operator between Hilbert spaces H -> H
% is continuous iff the operator norm is finite.

A linear operator between Hilbert spaces $U \to V$ is \textbf{continuous}
if and only if its operator norm is finite:
\[
    \sup_{\|u\|_U \le 1} \|L\,u\|_V \;<\; \infty.
\]

\begin{remark}
A \textbf{Hilbert space} is the natural infinite-dimensional
generalisation of Euclidean space $\mathbb{R}^d$: it is a
(possibly infinite-dimensional) vector space equipped with an
inner product that induces a complete norm.  Completeness means
that every Cauchy sequence converges, which is essential for
taking limits (e.g., of partial sums in Fourier series or of
approximating solutions to equations).
\end{remark}

\begin{remark}
The equivalence ``continuous $\Leftrightarrow$ bounded'' for
linear operators is a consequence of linearity: a linear map that
is continuous at one point is automatically continuous everywhere.
The \textbf{operator norm} $\|L\|_{\mathrm{op}} :=
\sup_{\|u\| \le 1} \|Lu\|$ measures the maximum ``stretching
factor'' of $L$.  The space of all bounded linear operators
$H \to H$, denoted $\mathcal{L}(H)$, is itself a Banach space
under this norm.
\end{remark}

%% =============================
%% Example: ell^2
%% =============================
\paragraph{Example: $\ell^2$.}

The space of square-summable sequences:
\[
    \ell^2
    \;=\;
    \left\{
        (u_n)_{n \in \mathbb{N}} \in \mathbb{R}^{\mathbb{N}}
        \;:\;
        \sum_{n=1}^{\infty} u_n^2 \;<\; \infty
    \right\}.
\]

The \textbf{inner product} on $\ell^2$ is:
\[
    \langle u, v \rangle
    \;=\;
    \sum_{k=1}^{\infty} u_k\,v_k.
\]

\begin{remark}
The space $\ell^2$ is the \textbf{prototypical separable Hilbert
space}.  By Parseval's identity, every separable Hilbert space
is isometrically isomorphic to $\ell^2$: choosing a CONS
$\{e_k\}$ and mapping $x \mapsto (\langle x, e_k\rangle)_k$
gives a bijective isometry $H \cong \ell^2$.  Thus, working in
$\ell^2$ is essentially the same as working in any other separable
Hilbert space.  The standard basis of $\ell^2$ is
$e_k = (0,\ldots,0,1,0,\ldots)$ with a $1$ in position $k$.
\end{remark}

%% =============================
%% Separability
%% =============================
\paragraph{Separability.}

A Hilbert space is called \textbf{separable} if there exists a
\textbf{countable dense subset} $H_0 \subset H$.

\begin{remark}
Separability is a mild assumption satisfied by all Hilbert spaces
that arise in practice (e.g., $L^2(M)$ for $\sigma$-finite
measure spaces $M$, and $\ell^2$).  Its key consequence is the
existence of a \textbf{countable} complete orthonormal system
(CONS) $\{e_k\}_{k=1}^\infty$, which allows us to represent
every element $x \in H$ as a convergent series
$x = \sum_k \langle x, e_k \rangle\, e_k$.  Without separability,
the CONS could be uncountable, and the theory of Gaussian
measures and trace-class operators becomes significantly more
delicate.
\end{remark}

%% =============================
%% Def 1.5: Gaussian measure
%% =============================
\subsection*{Definition 1.5: Gaussian measure}

\begin{definition}[Gaussian measure]
A probability measure $\mu$ on $(H,\,\mathcal{B}(H))$ is called
\textbf{Gaussian} if for all $h \in H$, the linear functional
\[
    \ell_h \colon H \;\longrightarrow\; \mathbb{R},
    \qquad
    g \;\longmapsto\; \langle g,\, h \rangle_H,
\]
is \textbf{normally distributed} under $\mu$. That is, there exist $m(h) \in \mathbb{R}$ and
$\sigma^2(h) \in \mathbb{R}^+$ such that for all $t \in \mathbb{R}$:
\[
    \phi_{\ell_h}(t) := \int e^{\,it\,\ell_h}\;d\mu
    \;=\;
    e^{\,it\,m(h) \;-\; \frac{1}{2}\,t^2\,\sigma^2(h)}.
\]
\end{definition}

\begin{remark}
Any linear functional on a Hilbert space can be represented as a
scalar product (inner product) with some element of $H$.
This is the \textbf{Riesz representation theorem}.
\end{remark}

\begin{remark}
Definition~1.5 says that a measure $\mu$ on $H$ is Gaussian if
\emph{every one-dimensional projection} of $\mu$ is Gaussian.
This is the infinite-dimensional analogue of the characterisation
of multivariate Gaussians in $\mathbb{R}^d$: a random vector
$X \in \mathbb{R}^d$ is Gaussian if and only if $a^T X$ is
(univariate) Gaussian for every $a \in \mathbb{R}^d$.
The inner product $\langle g, h\rangle_H$ plays the role of the
linear combination $a^T X$.
\end{remark}

% Characteristic function formulation:
% there exist m(h) in R and C^2(h) in R^+ such that
% for all t in R the characteristic function has the form:



% The dependence sigma^2(h) is quadratic in h.
% Linearity properties of m and sigma^2.

\paragraph{Linearity properties.}

Since $\ell_{h_1 + h_2} = \ell_{h_1} + \ell_{h_2}$
(the map $h \mapsto \ell_h$ is linear), we have:
\[
    \int \ell_{h_1 + h_2}\;d\mu
    \;=\;
    \int \ell_{h_1}\;d\mu
    \;+\;
    \int \ell_{h_2}\;d\mu,
\]
and therefore the \textbf{mean operator} is linear:
\[
    m(h_1 + h_2) \;=\; m(h_1) \;+\; m(h_2).
\]

\begin{remark}
The variance $\sigma^2(h)$ depends \textbf{quadratically} on $h$,
which leads to the covariance operator $Q$ satisfying
$\sigma^2(h) = \langle Qh,\, h \rangle$.
More precisely, the quadratic dependence means that $\sigma^2(\alpha h) = \alpha^2 \sigma^2(h)$ for all $\alpha \in \mathbb{R}$, and the map $(h_1, h_2) \mapsto \operatorname{Cov}(\ell_{h_1}, \ell_{h_2})$ is a \emph{bilinear form}.  The covariance operator $Q$ is the unique self-adjoint operator that represents this bilinear form via the inner product: $\operatorname{Cov}(\ell_{h_1}, \ell_{h_2}) = \langle Qh_1, h_2 \rangle$.
\end{remark}

%% =============================
%% Theorem 1.6: Characterisation of Gaussian measures
%% =============================
\subsection*{Theorem 1.6: Characterisation of Gaussian measures}

\begin{theorem}[Characterisation of Gaussian measures on Hilbert spaces]
A probability measure $\mu$ on $\bigl(H,\,\mathcal{B}(H)\bigr)$ is
\textbf{Gaussian} if and only if
\[
    \hat{\mu}(h) := \int_H e^{\,i\langle g,\, h \rangle_H}\;\mu(dg)
    \;=\;
    e^{\,i\langle m,\, h \rangle_H
      \;-\;\tfrac{1}{2}\,\langle Qh,\, h \rangle_H}
    \qquad \text{for all } h \in H,
\]
where:
\begin{enumerate}
    \item[\textup{a)}] $m \in H$ is the \textbf{mean vector},
    \item[\textup{b)}] $Q \colon H \to H$ is a linear, continuous operator on $H$ that is
    \begin{itemize}
        \item \textbf{symmetric} (self-adjoint), i.e.\
              $\langle Qh,\, g \rangle_H = \langle h,\, Qg \rangle_H$
              for all $h, g \in H$,
        \item \textbf{positive semi-definite}, i.e.\
              $\langle Qh,\, h \rangle_H \ge 0$ for all $h \in H$,
        \item of \textbf{finite trace}, i.e.\
              $\displaystyle
              \tr(Q)
              \;=\;
              \sum_{k=1}^{\infty}
              \langle Q e_k,\, e_k \rangle
              \;<\;\infty$
    \end{itemize}
    for any (hence for all) complete orthonormal system
    $\mathrm{(CONS)}$ $\{e_k\}$ of $H$.
\end{enumerate}
In this case we write $\mu = \mathcal{N}(m,\, Q)$, and $Q$ is called
the \textbf{covariance operator} of $\mu$.
\end{theorem}

\begin{remark}
The left-hand side of the characterisation formula is the
\textbf{characteristic function} (Fourier transform) of the measure $\mu$
evaluated at $h$.  In finite dimensions, a Gaussian
$\mathcal{N}(m, \Sigma)$ on $\mathbb{R}^d$ has characteristic function
$\exp\!\bigl(i\,m^T h - \tfrac{1}{2}\,h^T \Sigma\, h\bigr)$.
Theorem~1.6 is the infinite-dimensional generalisation: the covariance
\emph{matrix} $\Sigma$ is replaced by a covariance \emph{operator} $Q$,
and the finite-dimensional inner product $m^T h$ becomes
$\langle m, h\rangle_H$.
\end{remark}

\begin{remark}
The \textbf{trace-class} condition $\tr(Q) < \infty$ is essential and
has no analogue in finite dimensions (where every symmetric
positive semi-definite matrix automatically has finite trace).
In infinite dimensions, not every bounded symmetric positive operator
is trace-class.  The trace-class condition ensures that
$\mu$ is concentrated on $H$ (rather than on a larger space)
and that the Gaussian measure is well-defined as a \emph{Radon}
probability measure on $H$.
\end{remark}

\begin{remark}
The definition of $\tr(Q)$ is \textbf{independent of the choice of CONS}
$\{e_k\}$.  This follows from the fact that for a non-negative
self-adjoint operator $Q$, the quantity
$\sum_k \langle Q e_k, e_k\rangle$ is the same for every CONS
(a standard result in functional analysis).
\end{remark}

%% =============================
%% Properties of the Gaussian measure N(m,Q)
%% =============================
\paragraph{Properties of $\mathcal{N}(m,\,Q)$.}

In particular, the following hold:
\begin{enumerate}
    \item[\textup{(i)}]
    \textbf{Mean:}
    \[
        \int_H \langle x,\, h \rangle\;\mu(dx)
        \;=\;
        \langle m,\, h \rangle
        \qquad \text{for all } h \in H.
    \]

    \item[\textup{(ii)}]
    \textbf{Covariance:}
    \[
        \int_H
        \bigl(\langle x,\, h_1 \rangle - \langle m,\, h_1 \rangle\bigr)
        \bigl(\langle x,\, h_2 \rangle - \langle m,\, h_2 \rangle\bigr)
        \;\mu(dx)
        \;=\;
        \langle Q h_1,\, h_2 \rangle
        \;=\;
        \langle h_1,\, Q h_2 \rangle.
    \]

    \item[\textup{(iii)}]
    \textbf{Second moment / trace identity:}
    \begin{align*}
        \int_H \|x - m\|_H^2\;\mu(dx)
        &\;=\;
        \int_H \sum_{\ell=1}^{\infty}
        \langle x - m,\, e_\ell \rangle^2\;\mu(dx) \\[6pt]
        &\;=\;
        \sum_{\ell=1}^{\infty}
        \int_H \langle x - m,\, e_\ell \rangle^2\;\mu(dx) \\[6pt]
        &\;=\;
        \sum_{\ell=1}^{\infty}
        \langle Q e_\ell,\, e_\ell \rangle
        \;=\;
        \tr(Q).
    \end{align*}
\end{enumerate}

\begin{remark}
Property~(i) says that $m$ is the \textbf{Pettis integral}
(or weak integral) of the identity map under $\mu$: the mean
of each ``coordinate'' $\langle x, h\rangle$ is $\langle m, h\rangle$.

Property~(ii) shows that $Q$ encodes all second-order information.
Setting $h_1 = h_2 = h$ gives
$\mathrm{Var}_\mu\!\bigl(\langle x, h\rangle\bigr) = \langle Qh, h\rangle$,
recovering the variance formula from Definition~1.5.

Property~(iii) identifies the \textbf{expected squared norm} of the
centred random variable $x - m$ with the trace of the covariance
operator.  The interchange of sum and integral (second equality)
is justified by the \textbf{monotone convergence theorem},
since $\langle x - m, e_\ell\rangle^2 \ge 0$.
This gives a concrete probabilistic interpretation of $\tr(Q)$:
it is the total expected ``energy'' (squared norm) of the fluctuations
around the mean, distributed across all modes $e_\ell$ of the
Hilbert space.
\end{remark}

%% =============================
%% Parseval's identity
%% =============================
\paragraph{Parseval's identity.}

If $\{e_i\}$ is a CONS of $H$ and $x \in H$,
then writing $x_i := \langle x,\, e_i \rangle_H$ we have
\[
    \|x\|_H^2
    \;=\;
    \sum_{i=1}^{\infty} x_i^2
    \;=\;
    \sum_{i=1}^{\infty}
    \langle x,\, e_i \rangle^2.
\]

\begin{remark}
Parseval's identity says that every element of a separable Hilbert
space is completely determined by its Fourier coefficients
$x_i = \langle x, e_i\rangle$, and the norm is computed as
the $\ell^2$-norm of these coefficients.
In other words, the map $x \mapsto (\langle x, e_i\rangle)_{i\ge 1}$
is an \textbf{isometric isomorphism} from $H$ to $\ell^2$.
This identity was used in property~(iii) above when we expanded
$\|x - m\|_H^2 = \sum_\ell \langle x - m, e_\ell\rangle^2$.
\end{remark}

%% =============================
%% Example: Wiener measure on L^2
%% =============================
\subsection*{Example: Wiener process as a Gaussian measure on $L^2([0,T])$}

Let $(\Omega,\,\mathcal{F},\,P)$ be a probability space carrying a
\textbf{standard one-dimensional Brownian motion} $(B_t)_{t\in[0,T]}$
with continuous sample paths.

\paragraph{The sample-path mapping.}

Each realisation of the Brownian motion is a continuous function
on $[0,T]$.  We define the \textbf{sample-path mapping}
\[
    \beta \;\colon\; \Omega \;\longrightarrow\; C([0,T]),
    \qquad
    \omega \;\longmapsto\;
    \bigl(t \mapsto B_t(\omega)\bigr).
\]
Since continuous functions on a compact interval are square-integrable,
\[
    C([0,T]) \;\subseteq\; L^2([0,T])
    \qquad \text{(Hilbert space)},
\]
so every sample path $t \mapsto B_t(\omega)$ is an element of
$L^2([0,T])$.

\begin{remark}
The inclusion $C([0,T]) \subseteq L^2([0,T])$ is \emph{continuous}
but \emph{not} surjective: $L^2$ contains equivalence classes of
functions that may be highly irregular, whereas sample paths of
Brownian motion are a.s.\ continuous (hence in $C([0,T])$).
The important point here is that Brownian motion lives in a
``nice'' subset of the Hilbert space.
\end{remark}

\paragraph{Measurability and the induced measure.}

The map $\beta$ is measurable with respect to the Borel
$\sigma$-algebra $\mathcal{B}\!\bigl(L^2([0,T])\bigr)$.
We define the \textbf{Wiener measure} (or law of Brownian motion)
as the push-forward:
\[
    \mu(A)
    \;:=\;
    P\!\bigl(\beta \in A\bigr)
    \;=\;
    P\!\bigl(\{\omega \in \Omega : \beta(\omega) \in A\}\bigr),
    \qquad
    A \in \mathcal{B}\!\bigl(L^2([0,T])\bigr).
\]

Then $\mu$ is a \textbf{Gaussian measure} on $L^2([0,T])$ with
\textbf{mean zero} (i.e.\ $\mu$ is centred).

\begin{remark}
Saying that $\mu$ is a Gaussian measure on $L^2([0,T])$ means that
for every $h \in L^2([0,T])$, the real-valued random variable
$\langle \beta,\, h \rangle_{L^2} = \int_0^T B_s\,h(s)\, \mu(ds)$
is normally distributed.
This is indeed the case because $\int_0^T B_s\,h(s)\,\mu(ds)$ is a
Gaussian random variable for every deterministic $h \in L^2$
(as a limit of Riemann sums of jointly Gaussian random variables).
\end{remark}

\paragraph{Verification that the mean is zero.}

Let $\{e_k\}$ be any CONS of $L^2([0,T])$.
The $k$-th Fourier coefficient of a sample path is
\[
    x_k
    \;:=\;
    \langle B_\cdot,\, e_k \rangle_{L^2}
    \;=\;
    \int_0^T B_s\;e_k(s)\;ds.
\]
Its expectation is
\begin{align*}
    E(x_k)
    \;=\;
    E\!\left(\int_0^T B_s\;e_k(s)\;ds\right)
    &\;=\;
    \int_0^T E(B_s)\;e_k(s)\;ds
    \qquad \text{(Fubini)} \\[6pt]
    &\;=\;
    \int_0^T 0 \cdot e_k(s)\;ds
    \;=\; 0,
\end{align*}
since $E(B_s) = 0$ for a standard Brownian motion.
Because this holds for every $k$, the mean vector is $m = 0$.

\begin{remark}
The interchange of expectation and integral (application of
\textbf{Fubini's theorem}) is justified because
\[
    E\!\left(\int_0^T |B_s|\,|e_k(s)|\;ds\right)
    \;\le\;
    \int_0^T E(|B_s|)\,|e_k(s)|\;ds
    \;\le\;
    \int_0^T \sqrt{s}\;|e_k(s)|\;ds
    \;<\;\infty,
\]
using $E(|B_s|) = \sqrt{2s/\pi} \le \sqrt{s}$ and
$e_k \in L^2([0,T])$.
\end{remark}

\begin{remark}
One can further compute the \textbf{covariance operator} $Q$ of
the Wiener measure $\mu$.  For $h_1, h_2 \in L^2([0,T])$:
\[
    \langle Q h_1,\, h_2 \rangle_{L^2}
    \;=\;
    E\!\left(\int_0^T B_s\,h_1(s)\,ds
    \;\cdot\;
    \int_0^T B_r\,h_2(r)\,dr\right)
    \;=\;
    \int_0^T\!\!\int_0^T
    \min(s,r)\;h_1(s)\,h_2(r)\;ds\,dr,
\]
since $E(B_s B_r) = \min(s, r)$.  Thus $Q$ is the integral operator
with kernel $K(s,r) = \min(s,r)$:
\[
    (Qh)(t)
    \;=\;
    \int_0^T \min(t,s)\;h(s)\;ds.
\]
This operator is symmetric, positive semi-definite, and trace-class
(with $\tr(Q) = \int_0^T t\,dt = T^2/2$), confirming that $\mu$
satisfies all the conditions of Theorem~1.6.
\end{remark}

%% =============================
%% Covariance computation in detail
%% =============================
\paragraph{Detailed covariance computation.}

For $g, h \in L^2([0,T])$, the covariance of the Fourier
coefficients $x_g := \langle \beta, g\rangle_{L^2}$ and
$x_h := \langle \beta, h\rangle_{L^2}$ is computed as follows:

\begin{remark}
The following computation identifies the covariance operator $Q$ of the Wiener measure by directly computing $\operatorname{Cov}(\langle \beta, g\rangle, \langle \beta, h\rangle)$ and recognising the result as an inner product with an integral operator.  The key input is the covariance of Brownian motion, $E(B_s B_t) = \min(s,t)$, which encodes all the second-order information about the process.
\end{remark}
\begin{align*}
    \mathrm{Cov}(x_g,\, x_h)
    &\;=\;
    E\!\bigl(\langle \beta, g \rangle\,\langle \beta, h \rangle\bigr) \\[6pt]
    &\;=\;
    E\!\left(
        \int_0^T B_s\,g(s)\;ds
        \;\cdot\;
        \int_0^T B_t\,h(t)\;dt
    \right) \\[6pt]
    &\;=\;
    E\!\left(
        \int_0^T\!\!\int_0^T
        g(s)\,B_s\,B_t\,h(t)\;ds\;dt
    \right)
    \qquad \text{(double integral)} \\[6pt]
    &\;=\;
    \int_0^T\!\!\int_0^T
    g(s)\;E(B_s\,B_t)\;h(t)\;ds\;dt
    \qquad \text{(Fubini)} \\[6pt]
    &\;=\;
    \int_0^T\!\!\int_0^T
    g(s)\;(s \wedge t)\;h(t)\;ds\;dt,
\end{align*}
since $E(B_s B_t) = \min(s,t) = s \wedge t$.

Recognising the right-hand side as an $L^2$ inner product
with the integral operator, we obtain:
\[
    \mathrm{Cov}(x_g,\, x_h)
    \;=\;
    \int_0^T
    g(s)\;\underbrace{\left(\int_0^T (s \wedge t)\,h(t)\;dt\right)}_{=\,(Qh)(s)}\;ds
    \;=\;
    \langle g,\, Qh \rangle_{L^2},
\]
confirming that $Q$ is the covariance operator with
\[
    (Qh)(s)
    \;=\;
    \int_0^T \min(s,t)\;h(t)\;dt.
\]

%% =============================
%% Q as the inverse of the negative Laplacian
%% =============================
\paragraph{$Q$ as the inverse of a differential operator.}

We can split the integral defining $Qg$ by separating the regions
$s \le t$ and $s > t$:
\begin{align*}
    (Qg)(t)
    &\;=\;
    \int_0^T \min(s,t)\;g(s)\;ds \\[6pt]
    &\;=\;
    \int_0^t g(s)\,s\;ds
    \;+\;
    t\!\int_t^T g(s)\;ds.
\end{align*}

\begin{remark}
In the first integral ($0 \le s \le t$), we have $\min(s,t) = s$.
In the second integral ($t < s \le T$), we have $\min(s,t) = t$,
which is pulled out of the integral as a factor.
\end{remark}

Differentiating with respect to $t$ using the \textbf{Leibniz rule}:
\begin{align*}
    (Qg)'(t)
    &\;=\;
    \frac{d}{dt}\!\left[\int_0^t g(s)\,s\;ds\right]
    \;+\;
    \frac{d}{dt}\!\left[t\!\int_t^T g(s)\;ds\right] \\[6pt]
    &\;=\;
    g(t)\,t
    \;+\;
    \int_t^T g(s)\;ds
    \;-\; t\,g(t) \\[6pt]
    &\;=\;
    \int_t^T g(s)\;ds.
\end{align*}

Differentiating once more:
\[
    (Qg)''(t)
    \;=\;
    \frac{d}{dt}\!\left[\int_t^T g(s)\;ds\right]
    \;=\;
    -\,g(t).
\]

Therefore:
\[
    \boxed{-(Qg)''(t) \;=\; g(t),}
\]
which means $Q$ is the \textbf{inverse of the negative second derivative}
$-\partial_{tt}$, subject to the boundary conditions:
\begin{itemize}
    \item \textbf{Dirichlet at $t = 0$:}\quad
          $(Qg)(0)
          = \int_0^0 g(s)\,s\;ds + 0 \cdot \int_0^T g(s)\;ds
          = 0$.
    \item \textbf{Neumann at $t = T$:}\quad
          $(Qg)'(T)
          = \int_T^T g(s)\;ds
          = 0$.
\end{itemize}

That is:
\[
    Q \;=\; \left(-\frac{d^2}{dt^2}\right)^{\!-1}
    \quad
    \text{on } L^2([0,T]),
    \quad
    \text{with Dirichlet BC at } 0
    \text{ and Neumann BC at } T.
\]

\begin{remark}
This is a fundamental connection between probability and PDEs:
the covariance operator of the Wiener measure is the
\textbf{Green's operator} (inverse) of the negative Laplacian
$-\partial_{tt}$ on $[0,T]$ with mixed boundary conditions.
The kernel $K(s,t) = \min(s,t)$ is precisely the \textbf{Green's
function} of this boundary value problem.

Concretely, solving the ODE $-u''(t) = g(t)$ on $[0,T]$
with $u(0) = 0$ and $u'(T) = 0$ gives
$u(t) = \int_0^T \min(s,t)\,g(s)\,ds = (Qg)(t)$.
\end{remark}

\begin{remark}
The boundary conditions have a natural probabilistic interpretation:
\begin{itemize}
    \item The \textbf{Dirichlet condition} $(Qg)(0) = 0$ reflects
          the fact that Brownian motion starts at $B_0 = 0$
          (no randomness at time $0$).
    \item The \textbf{Neumann condition} $(Qg)'(T) = 0$ arises because
          Brownian motion has no ``preferred direction'' at the endpoint;
          its increments are symmetric.
\end{itemize}
\end{remark}

\begin{remark}
This relationship $Q = (-\partial_{tt})^{-1}$ is an instance
of a general principle: for many Gaussian processes, the covariance
operator is the inverse of a differential operator.
The \textbf{regularity} of the covariance kernel determines how
``smooth'' the typical sample paths are, and the order of the
differential operator reflects this.
Since $-\partial_{tt}$ is a second-order operator, its inverse $Q$
is a ``smoothing'' operator of order $2$, consistent with the fact
that Brownian paths are H\"older-$\tfrac{1}{2}^{-}$ continuous
(almost $\tfrac{1}{2}$ derivative of regularity).
\end{remark}

%% =============================
%% Why Q = (-Delta)^{-1} requires care
%% =============================
\subsection*{Why the Laplace operator is not usually invertible}

We have established that $Q = (-\Delta)^{-1}$, where $\Delta$ denotes
the \textbf{Laplace operator} (here $\Delta = \frac{d^2}{dt^2}$
in one dimension).

However, the Laplace operator is \textbf{not usually invertible}
on unbounded domains.  Consider $\Delta$ on $\mathbb{R}$ acting
on $L^2(\mathbb{R})$: the equation $-u''(t) = g(t)$ on $\mathbb{R}$
does not have a unique $L^2$ solution in general, because one can
always add linear functions (which solve the homogeneous equation
$u'' = 0$).  Moreover, the spectrum of $-\Delta$ on $L^2(\mathbb{R})$
is $[0, \infty)$ (purely continuous, no gap at zero), so $-\Delta$
has no bounded inverse.

\begin{remark}
To make $-\Delta$ invertible, one needs to impose
\textbf{boundary conditions} that pin down the solution.
In our example, the \textbf{Dirichlet condition} $u(0) = 0$
(corresponding to the fixed starting point $B_0 = 0$ of Brownian
motion) together with the Neumann condition at $T$ is precisely what
makes the operator invertible on the bounded interval $[0,T]$.
On a bounded domain with appropriate boundary conditions,
$-\Delta$ has a \textbf{discrete spectrum} $0 < \lambda_1 \le \lambda_2 \le \cdots$
and is invertible.
\end{remark}

%% =============================
%% Why Q = Id fails in infinite dimensions
%% =============================
\subsection*{Why $Q = \mathrm{Id}$ is not a valid covariance operator}

A natural first attempt might be to take the covariance operator
to be the \textbf{identity} $Q = \mathrm{Id}$
(i.e., every direction has variance $1$, as for a ``standard Gaussian''
in $\mathbb{R}^d$).  However, this fails in infinite dimensions.

\paragraph{Trace condition.}

Let $\{e_k\}_{k=1}^{\infty}$ be any CONS of $H$.  Then
\[
    \tr(\mathrm{Id})
    \;=\;
    \sum_{k=1}^{\infty}
    \langle e_k,\, e_k \rangle_H
    \;=\;
    \sum_{k=1}^{\infty} 1
    \;=\;
    +\infty.
\]
Since the trace is \textbf{not finite}, $Q = \mathrm{Id}$ is
\emph{not} trace-class, and by Theorem~1.6 there is \textbf{no}
Gaussian measure $\mathcal{N}(m,\,\mathrm{Id})$ on $H$.

\begin{remark}
This is a striking difference with finite dimensions.
In $\mathbb{R}^d$, the standard Gaussian $\mathcal{N}(0,\, I_d)$
is perfectly well-defined because
$\tr(I_d) = d < \infty$.
In an infinite-dimensional Hilbert space $H$, there are
\textbf{infinitely many independent directions}, each contributing
variance $1$, and the total variance
$E(\|X\|_H^2) = \tr(Q) = \sum_k 1 = \infty$
diverges.
Intuitively, the random vector would have infinite expected
squared norm and would ``not fit'' inside $H$.
\end{remark}

\paragraph{Geometric viewpoint: volumes of balls.}

There is a complementary geometric explanation.
In $\mathbb{R}^d$, the volume of the unit ball $B_1^d$ satisfies
\[
    \mathrm{Vol}(B_1^d)
    \;=\;
    \frac{\pi^{d/2}}{\Gamma(d/2 + 1)}
    \;\xrightarrow{d \to \infty}\;
    0.
\]
In infinite dimensions, interiors of balls have ``zero volume'':
the Gaussian mass escapes to higher and higher dimensions.
Formally, there is \textbf{no} $\sigma$-finite translation-invariant
measure (analogue of Lebesgue measure) on an infinite-dimensional
Hilbert space.

\begin{remark}
The fact that $\mathrm{Vol}(B_1^d) \to 0$ as $d \to \infty$
is a manifestation of the \textbf{curse of dimensionality}:
in high dimensions, most of the ``volume'' of a Gaussian
distribution is concentrated in a thin \textbf{annular shell}
far from the origin, not near the centre.
For a valid Gaussian measure on $H$, the covariance operator $Q$
must have \textbf{eigenvalues $\lambda_k \to 0$} fast enough
that $\sum_k \lambda_k = \tr(Q) < \infty$.
This forces the Gaussian to ``shrink'' in most directions,
compensating for having infinitely many of them.
\end{remark}

\begin{remark}
For the Wiener measure, the eigenvalues of $Q = (-\Delta)^{-1}$
on $[0,T]$ with Dirichlet--Neumann boundary conditions are
$\lambda_k = \frac{1}{\mu_k}$, where $\mu_k$ are the eigenvalues
of $-\Delta$.  Since $\mu_k \sim k^2$ for large $k$, the eigenvalues
$\lambda_k \sim k^{-2}$ decay to zero, and
$\tr(Q) = \sum_k \lambda_k < \infty$
(the series $\sum k^{-2}$ converges), confirming that the Wiener
measure is a well-defined Gaussian measure on $L^2([0,T])$.
\end{remark}

\end{document}
